\section{September 25, 2026}

\subsection{Invertible linear maps}

Write $\operatorname{Id}_V$ for the identity map on $V$, defined by
$\operatorname{Id}_V(v)=v$. As with matrix multiplication, we sometimes
write $ST$ for the composition $S\circ T$.

\begin{definition}[Invertibility]
  A linear map $T\in\mathcal L(V,W)$ is \emph{invertible} if there
  exists a linear map $S\in\mathcal L(W,V)$ such that
  \[
    S\circ T=\operatorname{Id}_V
    \qquad\text{and}\qquad
    T\circ S=\operatorname{Id}_W.
  \]
  Such a map $S$ is called an \emph{inverse} of $T$.
\end{definition}

\begin{proposition}[Uniqueness of the inverse]
  An invertible linear map has a unique inverse.
\end{proposition}

\begin{proof}
  Suppose $S_1,S_2\in\mathcal L(W,V)$ are both inverses of $T$.
  By associativity of composition,
  \begin{align*}
    S_1
      &=S_1\circ\operatorname{Id}_W \\
      &=S_1\circ(T\circ S_2) \\
      &=(S_1\circ T)\circ S_2 \\
      &=\operatorname{Id}_V\circ S_2=S_2.
  \end{align*}
\end{proof}

We denote the unique inverse by $T^{-1}\in\mathcal L(W,V)$. Thus
\[
  T^{-1}T=\operatorname{Id}_V
  \qquad\text{and}\qquad
  TT^{-1}=\operatorname{Id}_W.
\]

\begin{example}\label{ex:rotation-scaling-inverse}
  Define $T\colon\bb{R}^3\to\bb{R}^3$ by
  \[
    T\begin{bmatrix}x\\y\\z\end{bmatrix}
      =\begin{bmatrix}-y\\x\\4z\end{bmatrix}.
  \]
  Its values on the standard basis are
  \[
    T(e_1)=\begin{bmatrix}0\\1\\0\end{bmatrix},
    \qquad
    T(e_2)=\begin{bmatrix}-1\\0\\0\end{bmatrix},
    \qquad
    T(e_3)=\begin{bmatrix}0\\0\\4\end{bmatrix}.
  \]
  Hence its matrix with respect to the standard basis is
  \[
    A=\begin{bmatrix}0&-1&0\\1&0&0\\0&0&4\end{bmatrix}.
  \]
  To recover the input from an output $(x,y,z)$, reverse the first
  two coordinates with the appropriate sign and divide the third
  coordinate by $4$. This gives
  \[
    T^{-1}\begin{bmatrix}x\\y\\z\end{bmatrix}
      =\begin{bmatrix}y\\-x\\z/4\end{bmatrix},
    \qquad
    A^{-1}=\begin{bmatrix}0&1&0\\-1&0&0\\0&0&\tfrac14\end{bmatrix}.
  \]
  Indeed, direct multiplication gives $AA^{-1}=A^{-1}A=I_3$.
  Equivalently, augment $A$ by the identity matrix and perform the
  row operations $R_1\leftrightarrow R_2$, $R_2\leftarrow-R_2$,
  and $R_3\leftarrow\tfrac14 R_3$:
  \[
    \left[\begin{array}{ccc|ccc}
      0&-1&0&1&0&0\\
      1&0&0&0&1&0\\
      0&0&4&0&0&1
    \end{array}\right]
    \quad\longrightarrow\quad
    \left[\begin{array}{ccc|ccc}
      1&0&0&0&1&0\\
      0&1&0&-1&0&0\\
      0&0&1&0&0&\tfrac14
    \end{array}\right].
  \]
\end{example}

\begin{figure}[htbp]
  \centering
  \begin{tikzpicture}[>=Stealth, font=\small, scale=1.15]
    \draw[->, gray] (-1.8,0) -- (2.2,0) node[right] {$x$};
    \draw[->, gray] (0,-.7) -- (0,2.5) node[above] {$y$};
    \draw[->, very thick, blue!65!black] (0,0) -- (1.7,.85)
      node[right] {$(x,y)$};
    \draw[->, very thick, red!70!black] (0,0) -- (-.85,1.7)
      node[above left] {$(-y,x)$};
    \draw[->, thick] (26.565:1.3)
      arc[start angle=26.565, end angle=116.565, radius=1.3];
    \node at (.5,1.65) {$T$};
    \draw[->, thick, dashed] (116.565:2.15)
      arc[start angle=116.565, end angle=26.565, radius=2.15];
    \node at (1.15,2.1) {$T^{-1}$};
    \node[below left] at (0,0) {$0$};
  \end{tikzpicture}
  \caption{In the $xy$-plane, $T$ rotates by $90^\circ$
    counterclockwise and $T^{-1}$ rotates clockwise. In the third
    coordinate, $T$ multiplies by $4$ and $T^{-1}$ divides by $4$.}
  \label{fig:rotation-inverse}
\end{figure}

\subsection{Invertibility and bijectivity}

\begin{proposition}\label{prop:invertible-iff-bijective}
  A linear map is invertible if and only if it is bijective.
\end{proposition}

\begin{proof}
  Let $T\in\mathcal L(V,W)$. Suppose first that $T$ is invertible.
  If $T(u)=T(v)$, applying $T^{-1}$ gives
  \[
    u=T^{-1}(T(u))=T^{-1}(T(v))=v.
  \]
  Thus $T$ is injective. For every $w\in W$, the vector
  $v=T^{-1}(w)$ satisfies
  \[
    T(v)=T(T^{-1}(w))=w,
  \]
  so $T$ is also surjective.

  Conversely, suppose that $T$ is bijective. For each $w\in W$,
  surjectivity provides $v\in V$ with $T(v)=w$, and injectivity
  makes this vector unique. Define $S\colon W\to V$ by assigning
  to $w$ this unique vector. By construction,
  $T\circ S=\operatorname{Id}_W$. Also, for $v\in V$,
  \[
    T(S(T(v)))=T(v),
  \]
  so injectivity gives $S(T(v))=v$. Thus
  $S\circ T=\operatorname{Id}_V$.

  It remains to check that $S$ is linear. For $w_1,w_2\in W$,
  \begin{align*}
    T(S(w_1)+S(w_2))
      &=T(S(w_1))+T(S(w_2)) \\
      &=w_1+w_2 \\
      &=T(S(w_1+w_2)).
  \end{align*}
  Injectivity of $T$ yields
  \[
    S(w_1+w_2)=S(w_1)+S(w_2).
  \]
  Similarly, for $\lambda\in\bb{F}$ and $w\in W$,
  \[
    T(\lambda S(w))=\lambda T(S(w))=\lambda w=T(S(\lambda w)),
  \]
  so $S(\lambda w)=\lambda S(w)$. Therefore
  $S\in\mathcal L(W,V)$ and $S$ is an inverse of $T$.
\end{proof}

\subsection{Why the dimension hypotheses matter}

In general, injectivity alone or surjectivity alone does not imply
invertibility.

\begin{example}
  On the space $\mathcal P(\bb{R})$ of all real polynomials, define
  \[
    M\colon\mathcal P(\bb{R})\to\mathcal P(\bb{R}),
    \qquad (Mp)(x)=x^2p(x).
  \]
  This map is linear and injective: if $x^2p(x)$ is the zero
  polynomial, then $p$ is the zero polynomial. It is not surjective,
  because the constant polynomial $1$ does not belong to its image.
  Thus $M$ is not invertible, even though its domain and codomain
  are the same vector space.
\end{example}

\begin{example}[A complementary example]
  Differentiation
  \[
    D\colon\mathcal P(\bb{R})\to\mathcal P(\bb{R}),
    \qquad D(p)=p',
  \]
  is surjective, since every real polynomial has a polynomial
  antiderivative. It is not injective, since every constant polynomial
  belongs to $\ker D$. Therefore $D$ is not invertible.
\end{example}

\begin{proposition}\label{prop:equal-dimension-invertibility}
  Suppose that $V$ and $W$ are finite-dimensional with
  $\dim V=\dim W$, and let $T\in\mathcal L(V,W)$. Then the
  following are equivalent:
  \begin{enumerate}[label=(\roman*)]
    \item $T$ is invertible;
    \item $T$ is injective;
    \item $T$ is surjective.
  \end{enumerate}
\end{proposition}

\begin{proof}
  By rank--nullity and the equality of dimensions,
  \begin{align*}
    T\text{ is injective}
      &\quad\Longleftrightarrow\quad \ker T=\{0\} \\
      &\quad\Longleftrightarrow\quad \dim\ker T=0 \\
      &\quad\Longleftrightarrow\quad
        \dim\operatorname{im}T=\dim V=\dim W \\
      &\quad\Longleftrightarrow\quad \operatorname{im}T=W \\
      &\quad\Longleftrightarrow\quad T\text{ is surjective}.
  \end{align*}
  Here we used the fact that a subspace of a finite-dimensional space
  with the same dimension is the whole space. Thus either
  injectivity or surjectivity implies bijectivity, which is equivalent
  to invertibility by Proposition~\ref{prop:invertible-iff-bijective}.
\end{proof}

\subsection{One-sided inverses in equal dimensions}

\begin{proposition}\label{prop:one-sided-inverse}
  Suppose that $V$ and $W$ are finite-dimensional with
  $\dim V=\dim W$. If $T\in\mathcal L(V,W)$ and
  $S\in\mathcal L(W,V)$, then
  \[
    ST=\operatorname{Id}_V
    \quad\Longleftrightarrow\quad
    TS=\operatorname{Id}_W.
  \]
  In either case, $S=T^{-1}$.
\end{proposition}

\begin{proof}
  Suppose $ST=\operatorname{Id}_V$. If $v\in\ker T$, then
  \[
    v=(ST)(v)=S(T(v))=S(0)=0.
  \]
  Hence $T$ is injective, and
  Proposition~\ref{prop:equal-dimension-invertibility} makes it
  invertible. Consequently,
  \[
    S=S(TT^{-1})=(ST)T^{-1}=T^{-1},
  \]
  which implies $TS=\operatorname{Id}_W$.

  Conversely, suppose $TS=\operatorname{Id}_W$. For every $w\in W$,
  \[
    w=T(S(w)),
  \]
  so $T$ is surjective and hence invertible by the same proposition.
  Therefore
  \[
    S=(T^{-1}T)S=T^{-1}(TS)=T^{-1},
  \]
  and $ST=\operatorname{Id}_V$.
\end{proof}
